\documentclass[12pt,a4paper]{article}
\usepackage[utf8]{inputenc}
\usepackage{ctex} % 支持中文
\usepackage{amsmath, amssymb, amsthm}
\usepackage{graphicx}
\usepackage{listings}
\usepackage{xcolor}
\usepackage{tabularx}
\usepackage{booktabs}
\usepackage{float} % Required for the H float option

\lstset{
    language=Python, % 添加语言支持
    basicstyle=\ttfamily\small,
    numbers=left,
    numberstyle=\tiny,
    keywordstyle=\color{blue},
    commentstyle=\color{gray},
    stringstyle=\color{red},
    breaklines=true,
    frame=single,
    captionpos=b
}

\title{计算流体力学作业4}
\author{郑恒2200011086}
\date{\today}

\begin{document}

\maketitle

\section{问题介绍}
本次作业研究一个矩形平板的稳态温度场分布问题。平板尺寸为 $15 \, \text{cm} \times 12 \, \text{cm}$，假设板的表面是绝热的，仅在四条边有热流通过。一条 $15 \, \text{cm}$ 的边温度保持为 $100^\circ \text{C}$，其余三条边温度为 $20^\circ \text{C}$。

为求解该问题，采用离散代数方法，使用SOR（超松弛高斯赛德尔迭代）方法计算温度场，并绘制等温线。具体步骤如下：
\begin{enumerate}
    \item 将矩形平板离散为网格点，建立离散方程；
    \item 选择不同的松弛因子 $\omega$，比较收敛速度；
    \item 采用不同的网格尺寸，观察最佳松弛因子的变化；
    \item 使用Python编程实现计算，并绘制温度场的等温线图。
\end{enumerate}
\newpage
\section{代码实现}


\newpage   
\section{AI工具使用说明表}
\begin{table}[!htbp]
    \centering
    \begin{tabular}{|c|c|c|}
        \hline
        \textbf{AI名称} & \textbf{生成代码功能} & \textbf{使用内容} \\
        \hline
        Copilot & latex格式框架 & figure参数调整、图片插入\\
        \hline
        Deepseek & python绘图调整 & 90-104、156-166、203-210、234-241行图绘制的具体参数调整\\
        \hline
        Deepseek & gitignore文件忽略 & 全由ai添加\\
        \hline
\end{tabular}
\end{table}

commit 记录如下图：
\begin{figure}[H]
    \centering
    \includegraphics[width=0.8\textwidth]{commits.pdf}
    \caption{commit 记录}
    \label{fig:commit_record}
\end{figure}

\end{document}